\section{Muestreo de la posterior difusiva: DPS}

Cuando estamos dispuestos a aceptar una aproximación más gruesa a cambio de una aplicabilidad más amplia, obtenemos el método \textbf{DPS (Diffusion Posterior Sampling)}.

\subsection{Aproximación Delta más gruesa}

DPS utiliza una aproximación más extrema que la gaussiana: aproxima directamente $p(x_0 | x_t)$ como una distribución Delta:

$$
p(x_0 | x_t) \approx \delta(x_0 - \hat{x}_0)
$$

Donde $\hat{x}_0$ es la estimación de Tweedie. Esto significa que asumimos que $x_0$ puede recuperarse determinísticamente a partir de $x_t$.

\begin{warningbox}{Costos y beneficios de la aproximación Delta}
\textbf{Costo}: Es una aproximación peor que la gaussiana. Si $p(x_0|x_t)$ realmente fuera una distribución Delta, no sería necesario realizar un proceso iterativo de eliminación del ruido.

\textbf{Beneficio}: El cálculo se vuelve extremadamente simple y \textbf{ya no es necesario que $\mathcal{A}$ sea lineal} —se puede utilizar cualquier mapa hacia adelante diferenciable.
\end{warningbox}

\subsection{Puntuación de verosimilitud del DPS}

Bajo la aproximación Delta:

$$
p(y | x_t) \approx p(y | x_0 = \hat{x}_0) = \mathcal{N}(y; \mathcal{A}(\hat{x}_0), \sigma_y^2 I)
$$

\begin{importantbox}{Fórmula de guía del DPS}
La puntuación de verosimilitud del DPS es:
$$
\nabla_{x_t} \log p(y | x_t) \approx -\frac{1}{\sigma_y^2} \nabla_{x_t} \left\| y - \mathcal{A}(\hat{x}_0) \right\|_2^2
$$
Donde $\hat{x}_0 = \hat{x}_0(x_t)$ es la estimación de eliminación del ruido calculada a partir de $x_t$ mediante la fórmula de Tweedie, y $\nabla_{x_t}$ representa el gradiente respecto a $x_t$ (que requiere retropropagación a través de la dependencia de $\hat{x}_0$ en $x_t$).

\textbf{Significado intuitivo}: Comparar la diferencia entre el resultado del mapa hacia adelante aplicado a la estimación de Tweedie, $\mathcal{A}(\hat{x}_0)$, y la observación real $y$, calcular el gradiente de la pérdida L2 respecto a $x_t$ para servir como señal de guía.
\end{importantbox}

\subsection{Flujo del algoritmo DPS}

La implementación del DPS es extremadamente concisa: solo se requiere agregar una línea de código al ciclo de muestreo estándar de DDPM:

\begin{lstlisting}[caption={Pseudocódigo de guía DPS}]
# Bucle estándar de muestreo DDPM
for t in reversed(range(T)):
    # 1. Paso estándar de eliminación del ruido
    eps_pred = model(x_t, t)           # predecir el ruido
    x0_hat = tweedie_estimate(x_t, eps_pred, t)  # Tweedie
    x_t_minus_1 = ddpm_step(x_t, eps_pred, t)    # paso estándar

    # 2. Guía DPS (¡solo una línea adicional!)
    loss = torch.norm(y - A(x0_hat)) ** 2
    grad = torch.autograd.grad(loss, x_t)[0]
    x_t_minus_1 = x_t_minus_1 - step_size * grad
\end{lstlisting}

\begin{importantbox}{Dos implementaciones equivalentes del DPS}
Bajo el marco de DDPM, las siguientes dos implementaciones son equivalentes:
\begin{enumerate}
    \item Modificar la predicción de ruido $\epsilon$ y luego ejecutar el paso de eliminación del ruido.
    \item Ejecutar primero el paso estándar de eliminación del ruido para obtener $x_{t-1}'$, y luego agregar directamente el término de guía en $x_{t-1}'$.
\end{enumerate}
Esto se debe a que $x_{t-1}$ es una función afín de $\epsilon$, por lo que ambas formas son matemáticamente completamente equivalentes. La segunda implementación es más intuitiva y concisa.
\end{importantbox}

\subsection{Alcance del DPS}

\begin{knowledgebox}{El DPS soporta mapas hacia adelante no lineales}
Una ventaja clave del DPS es que \textbf{no requiere} que $\mathcal{A}$ sea lineal. Siempre que $\mathcal{A}$ sea diferenciable (para permitir la retropropagación y el cálculo de gradientes), se puede aplicar el DPS. Esto permite al DPS manejar una gama más amplia de problemas inversos, tales como:
\begin{itemize}
    \item Recuperación de imágenes con distorsión no lineal
    \item Reconstrucción 3D que involucra renderizadores basados en redes neuronales
    \item Recuperación de señales basada en modelos físicos no lineales
\end{itemize}
El costo es una precisión de aproximación menor (aproximación Delta), pero en la práctica todavía produce resultados muy buenos.
\end{knowledgebox}

\subsection{Resumen de la sección}

DPS convierte el cálculo de la puntuación de verosimilitud en un gradiente de una pérdida L2 simple al adoptar la aproximación Delta más sencilla. Su implementación solo requiere agregar un código de actualización de gradiente adicional al ciclo de muestreo estándar. Aunque la aproximación es más gruesa, esto a cambio soporta mapas hacia adelante no lineales, ampliando significativamente el alcance de aplicación.
